\chapter{Concluzii}

În această lucrare am evaluat trei sisteme SQL distribuite, și anume: CockroachDB, TiDB și Citus pe suita analitică TPC-H la factor de scară 30, pe configurații de cluster cu 3, 5, 10 și 20 de noduri. Analiza celor 22 de interogări împreună cu interpretarea planurilor de execuție corespunzătoare a evidențiat trei arhitecturi distincte, fiecare cu aplicabilitatea ei.

Concluzia generală a studiului este că alegerea sistemului optim pentru o aplicație analitică distribuită depinde de combinația dintre profilul dominant al cerințelor de lucru și nivelul de acoperire SQL necesar, nu doar de performanța brută. Pentru cerințe de tip analitic OLAP și control complet asupra strategiei de partiționare Citus reprezintă alegerea cu cele mai bune rezultate. Pentru cerințe mixte sau pentru cazuri în care acoperirea SQL și predictibiltatea primează asupra latenței absolute, CockroachDB constituie o alegere echilibrată, rămânând singurul sistem care a rulat toate interogările. Pentru cerințe mixte, TiDB rămâne competitiv prin mecanismele sale de separare.

\section{Direcții viitoare de cercetare}

Sunt mai multe direcții de continuare a lucrării. O primă direcție este reevaluarea TiDB prin utilizarea componentei TiFlash, un motor columnar replicat care nu a făcut parte din configurația acestui studiu dar care ar fi putut modifica semnificativ rezultate. O altă direcție de dezvoltare al acestui studiu este extinderea evaluării pe suita TPC-C, care ar permite caracterizarea sistemelor pe profil tranzacțional OLTP și compararea echilibrului, aceste sisteme fiind proiectate în principal pentru acest tip de procesare. O a treia direcție este compararea cu sisteme analitice native columnare. O a patra direcție ar fi efectul indecșilor creați manuali fie la recomandarea motoarelor, fie prin analiza manuală, sistemul CRDB sugerând explicit indecși pentru fiecare interogare. O direcție complementară este evaluarea sub concurență, prin execuția mai multor interogări simultan.  O altă direcție firească este creșterea factorului de scară pentru a evalua comportamentul sistemelor la volume mai mari. 
